%!TEX root = trajectory-grouping.tex
% -- Author: Jannes Bantje, j.bantje@wwu.de
\documentclass[a4paper,index=totoc,toc=bibliography,fontsize=11pt,DIV=13,headinclude,BCOR=5mm,cleardoublepage=empty,draft]{scrreprt}
\usepackage{scrpage2} % wie fancyhdr, nur optimiert auf KOMA-Skript, leicht andere Syntax
\usepackage{etoolbox,letltxmacro} % zum Programmieren
\usepackage[final]{graphicx}
\robustify\rotatebox

% -- Farben
\usepackage[x11names]{xcolor}
\definecolor{dark_gray}{gray}{0.45}
\definecolor{light_gray}{gray}{0.6}
\definecolor{fb10_blue}{cmyk}{0.8,0.4,0.13,0.07}

\usepackage[utf8]{inputenc}
\usepackage[T1]{fontenc}
\usepackage{mathtools} % beinhaltet amsmath
\mathtoolsset{centercolon,showmanualtags}
\newtagform{brackets}{[}{]}
\usetagform{brackets}
\usepackage{fix-cm}
\usepackage[bbgreekl]{mathbbol}
\usepackage{amssymb} % zusätzliche Symbole
\usepackage{nicefrac} % schräge Brüche
\usepackage{faktor}
\newcommand{\Faktor}[1]{\faktor[\textstyle]{#1}}
\usepackage{xfrac}
\usepackage{cancel}
\usepackage{mathdots} % Verbesserung von Punkten wie zB \ldots
\usepackage[bb=px]{mathalfa}
\usepackage{centernot}

\usepackage{silence}
\WarningFilter{latexfont}{Size substitutions with differences}
\WarningFilter{latexfont}{Font shape `U/bbold/m/n' in size}

%!TEX root = trajectory-grouping.tex

% -- Zum Finetuning von Befehlen
\makeatletter
\newcommand{\raisemath}[1]{\mathpalette{\raisem@th{#1}}}
\newcommand{\raisem@th}[3]{\raisebox{#1}{$#2#3$}}
\makeatother

% -- Box mit vorgegebener minimaler Länge
\DeclareRobustCommand{\minwidthbox}[2]{%
  \ifmmode
    \expandafter\mathmakebox
  \else
    \expandafter\makebox
  \fi
  [\ifdim#2<\width\width\else#2\fi]{#1}%
}

% -- farbiges Untersteichen im Mathe-Modus
\def\mathul#1#2{\color{#1}\underline{{\color{black}#2}}\color{black}} 


% -- besserer underbrace befehl
\newcommand{\Underbrace}[2]{{\underbrace{#1}_{#2}}}





%--Mengen
\newcommand\SetSymbol[1][]{\nonscript\:#1\vert\allowbreak\nonscript\:\mathopen{}}
\providecommand\given{} % to make it exist
\DeclarePairedDelimiterX\set[1]\{\}{\renewcommand\given{\SetSymbol[\delimsize]}#1}

% -- Betrag und Skalarprodukt
\DeclarePairedDelimiter{\abs}{\lvert}{\rvert}
\DeclarePairedDelimiterX\skal[2]{\langle}{\rangle}{#1\,,\,#2}

%--Umklammern
\DeclarePairedDelimiter\enbrace{(}{)}
\DeclarePairedDelimiter\benbrace{[}{]}
\DeclarePairedDelimiter\homo{\llbracket}{\rrbracket}
\newcommand{\ssbrace}[1]{{\scriptscriptstyle\enbrace{#1}}}

%--Norm
\DeclarePairedDelimiter\norm{\Vert}{\Vert}






% - verbessertes Integral
\newcommand{\Int}[1]{\int_{\mathrlap{#1}}\,\,}

% -- Alles spezielles aus der Differentialtopologie
\newcommand{\Tang}{\ensuremath{\mathrm{T}\mkern-0.85mu}}
\newcommand{\mathd}{\ensuremath{\mathrm{d}\mkern-0.7mu}}
\newcommand{\diff}[2]{\ensuremath{\frac{{\partial #1}}{{\partial #2}}}}
\newcommand{\diffs}[2]{\ensuremath{\partial #1/\partial #2}}
\newcommand{\diffd}[2]{\ensuremath{\frac{\mathd #1}{\mathd #2} }}
\DeclareMathOperator{\Lie}{L}
\DeclareMathOperator{\Ad}{Ad}  
\DeclareMathOperator{\ad}{ad}
\DeclareMathOperator{\Spin}{Spin}
% \DeclareMathOperator{\spin}{spin}
\newcommand{\spin}{\mathop{\mathfrak{spin}}}
\newcommand{\Ce}{\mathcal{C}}
\DeclareMathOperator{\vol}{vol}



%--Abbildungsdefinition
\newcommand{\mapdef}[5]{%
	\[
		\begin{array}{rcl}
			\textstyle #1 &\xrightarrow{\minwidthbox{#5}{2em}} & \textstyle #2 \\[0.5ex]
			\textstyle #3 &\xmapsto{\minwidthbox{\mbox{ }}{2em}} & \textstyle #4
		\end{array}
	\]
}

%--modifiziertes Stackrel
\newcommand{\StackText}[2]{\stackrel{\mbox{\scriptsize #1}}{#2}}
\newcommand{\StackTextClap}[2]{\stackrel{\mathclap{\mbox{\scriptsize #1}}}{#2}}



 % Liste mit Mathebefehlen laden

% -- TikZ-Kram
\usepackage{tikz-cd}
\usetikzlibrary{quotes,babel,arrows.meta}
\tikzset{>=Straight Barb[]}
\tikzcdset{arrow style=math font}

% -- tikz/externalize
\usetikzlibrary{external}
\tikzexternalize[prefix=Bilder/tikz/, up to date check=diff]
\tikzset{external/system call={xelatex \tikzexternalcheckshellescape %-- verwende LuaLaTeX, wegen dynamischer Speicherallokation
    -halt-on-error -interaction=batchmode -jobname "\image" "\texsource"}}
\AtBeginEnvironment{tikzcd}{\tikzexternaldisable}
\AtEndEnvironment{tikzcd}{\tikzexternalenable}
\pgfkeys{/pgf/images/include external/.code=\includegraphics{#1}}

% -- Einstellungen bezüglich Schriftart etc.
\usepackage{ccfonts}
\usepackage[semibold]{sourcesanspro}
\usepackage[semibold,scale=.95]{sourcecodepro}
\usepackage{eulervm}

% \addtokomafont{disposition}{\sffamily}

\usepackage{fontawesome}
\usepackage[final]{microtype}
\flushbottom
\KOMAoptions{DIV=last}


% -- Einstellungen bezüglich der Sprache, Silbentrennung etc.
\usepackage[ngerman]{babel}
% \usepackage{polyglossia}
% \setmainlanguage[spelling=new,babelshorthands=true]{german}
% \setotherlanguage[variant=british]{english}

% -- BibLaTeX 
\usepackage[%
	backend=biber,
	sortlocale=auto,
	natbib,
	hyperref,
	backref,
	style=alphabetic
	]%
{biblatex}
\renewcommand*{\mkbibnamelast}[1]{%
  \ifmknamesc{\textsc{#1}}{#1}}

\renewcommand*{\mkbibnameprefix}[1]{%
  \ifboolexpr{ test {\ifmknamesc} and test {\ifuseprefix} }
    {\textsc{#1}}
    {#1}}

\def\ifmknamesc{%
  \ifboolexpr{ test {\ifcurrentname{labelname}}
               or test {\ifcurrentname{author}}
               or ( test {\ifnameundef{author}} and test {\ifcurrentname{editor}} ) }}
\addbibresource{quellen.bib}

% -- Konfiguration von Hyperref
\usepackage[hidelinks, pdfpagelabels, bookmarksopen=true, bookmarksnumbered=true, linkcolor=black, urlcolor=SkyBlue2, plainpages=false,pagebackref, citecolor=black, hypertexnames=true, pdfauthor={Jannes Bantje}, pdfborderstyle={/S/U}, linkbordercolor=SkyBlue2, colorlinks=false,backref=false]{hyperref}
\hypersetup{final}

\usepackage[nameinlink,noabbrev]{cleveref}

% -- \url aufhübschen
\newcommand{\appendLink}[1]{#1\,\faExternalLink}
\newcommand{\hrefsym}[2]{\href{#1}{\texttt{\appendLink{#2}}}}
\renewcommand{\url}[1]{\hrefsym{#1}{\nolinkurl{#1}}}

% -- Aufzählungen, Anführungszeichen etc.
\usepackage[shortlabels,inline]{enumitem}
\setlist[itemize,1]{label=\scriptsize$\blacksquare$}
\usepackage[autostyle,german=quotes]{csquotes}


% -- Indexverarbeitung
\usepackage{makeidx}
\newcommand{\bet}[1]{\emph{#1}}
\newcommand{\Index}[1]{\bet{#1}\index{#1}}
\makeindex
% \setindexpreamble{{\noindent \itshape Die Seitenzahlen sind mit Hyperlinks zu den entsprechenden Seiten versehen, also anklickbar}\par \bigskip}
\renewcommand{\indexpagestyle}{scrheadings}

% -- Alles zum Thema Gliederung
% \usepackage{chngcntr}
% \counterwithin*{equation}{section}
\numberwithin{equation}{section}
% \setcounter{tocdepth}{3}

% -- theorem packages
\let\openbox\relax
\usepackage{amsthm}
\usepackage{thmtools,thm-restate}
\usepackage{mdframed}
\renewcommand{\listtheoremname}{Übersicht aller Aussagen}

% -- Theoreme als PDF-Lesezeichen
\usepackage{bookmark}
\bookmarksetup{open,numbered}
\makeatletter
\newcommand*{\theorembookmark}{%
  \bookmark[
    dest=\@currentHref,
    rellevel=1,
    keeplevel,
  ]{%
    \thmt@thmname\space\csname the\thmt@envname\endcsname
    \ifx\thmt@shortoptarg\@empty
    \else
      \space(\thmt@shortoptarg)%
    \fi
  }%
}   
\makeatother

% -- Definition der einzelnen Umgebungen
\declaretheoremstyle[%
	headfont=\sffamily\bfseries,
	notefont=\normalfont\sffamily,
	bodyfont=\normalfont,
	headformat=\NUMBER\ \NAME\NOTE,
	headpunct=.,
	postheadspace=1em,
	spaceabove=15pt,spacebelow=10pt,
	shaded={bgcolor=gray!20},
	% mdframed={%
	% 	backgroundcolor=gray!20,
	% 	innertopmargin=0pt,
	% 	innerleftmargin=0pt,
	% 	% roundcorner=5pt,
	% 	innerbottommargin=0pt,
	% 	splittopskip = 0pt,
	% 	% skipbelow=6pt,
	% 	% skipbelow=6pt,
	% 	topline=false,bottomline=false,leftline=false,rightline=false
	% },
	postheadhook=\theorembookmark]%
{mainstyle}
\declaretheoremstyle[%
	headfont=\sffamily\bfseries,
	notefont=\normalfont\sffamily,
	bodyfont=\normalfont,
	headformat=\NUMBER\ \NAME\NOTE,
	headpunct=.,
	postheadspace=1em,
	spaceabove=15pt,spacebelow=10pt,
	shaded={bgcolor=fb10_blue!20},
	postheadhook=\theorembookmark]%
{mainstyle_blue}
\declaretheoremstyle[%
	headfont=\sffamily\bfseries,
	notefont=\normalfont\sffamily,
	bodyfont=\normalfont,
	headformat=\NUMBER\ \NAME\NOTE,
	headpunct=.,
	postheadspace=1em,
	spaceabove=15pt,spacebelow=10pt,
	postheadhook=\theorembookmark]%
{mainstyle_unshaded}
\declaretheoremstyle[%
	headfont=\sffamily\bfseries,
	notefont=\normalfont\sffamily,
	bodyfont=\normalfont,
	headformat=\NUMBER\NAME\NOTE,
	headpunct=.,
	postheadspace=1em,
	spaceabove=15pt,spacebelow=10pt,
	% shaded={bgcolor=gray!20},
	postheadhook=\theorembookmark]%
{mainstyle_unnumbered}
\declaretheoremstyle[%
	headfont=\sffamily\bfseries,
	notefont=\normalfont\sffamily,
	bodyfont=\normalfont,
	headformat=swapnumber,
	headpunct=.,
	postheadspace=1em,
	spaceabove=15pt,spacebelow=10pt,
	shaded={bgcolor=gray!20},
	postheadhook=\theorembookmark,
	qed=\qedsymbol]%
{mainstyleB}
\declaretheorem[name=Definition,parent=section,style=mainstyle_blue]{definition}
\declaretheorem[name=Definition \& Proposition,refname=Proposition,sharenumber=definition,style=mainstyle_blue]{definitionP}
\declaretheorem[name=Definition,numbered=no,style=mainstyle_unnumbered]{definition*}
\declaretheorem[name=Theorem,sharenumber=definition,style=mainstyle]{theorem}
\declaretheorem[name=Theorem,numbered=no,style=mainstyle_unnumbered]{theorem*}
\declaretheorem[name=Proposition,sharenumber=definition,style=mainstyle,refname=Proposition]{proposition}
\declaretheorem[name=Lemma,sharenumber=definition,style=mainstyle]{lemma}
\declaretheorem[name=Satz,refname=Satz,sharenumber=definition,style=mainstyle]{satz}
\declaretheorem[name=Satz,sharenumber=definition,style=mainstyle_unshaded]{satzUnshaded}
\declaretheorem[name=Definition,sharenumber=definition,style=mainstyle_unshaded]{definitionUnshaded}
\declaretheorem[name=Satz,numbered=no,style=mainstyle_unnumbered]{satz*}
\declaretheorem[name=Korollar,sharenumber=definition,style=mainstyle]{korollar}
\declaretheorem[name=Korollar,sharenumber=definition,style=mainstyleB]{korollarB}
\declaretheorem[name=Frage,numbered=no,style=mainstyle_unnumbered]{frage}
\declaretheorem[name=Frage,sharenumber=definition,style=mainstyle_unshaded]{frageA}
\declaretheorem[name=Erinnerung,sharenumber=definition,style=mainstyle_unshaded]{erinnerungA}
\declaretheorem[name=Ausblick,sharenumber=definition,style=mainstyle_unshaded]{ausblick}
\declaretheorem[name=Konvention,sharenumber=definition,style=mainstyle]{konvention}
\declaretheorem[name=Notation,sharenumber=definition,style=mainstyle_unshaded]{notation}
\declaretheorem[name=Bemerkung,sharenumber=definition,style=mainstyle_unshaded]{bemerkung}
\declaretheorem[name=Bemerkung,numbered=no,style=mainstyle_unnumbered]{bemerkung*}
\declaretheorem[name=Beispiel,sharenumber=definition,style=mainstyle_unshaded]{beispiel}
\declaretheorem[name=Beispiel,numbered=no,style=mainstyle_unnumbered]{beispiel*}
\declaretheorem[name=Exkurs,numbered=no,style=mainstyle_unnumbered]{exkurs*}

% english versions
\declaretheorem[name=Remark,sharenumber=definition,style=mainstyle_unshaded]{remark}
\declaretheorem[name=Remark,numbered=no,style=mainstyle_unnumbered]{remark*}
\declaretheorem[name=Example,sharenumber=definition,style=mainstyle_unshaded]{example}
\declaretheorem[name=Corollary,sharenumber=definition,style=mainstyle]{corollary}

% -- Beweise
\declaretheoremstyle[headfont=\bfseries\scshape,bodyfont=\normalfont,headpunct=:,postheadspace=1em,spacebelow=12pt,spaceabove=2pt,qed=\qedsymbol]{beweise}
\declaretheoremstyle[headfont=\bfseries\scshape,bodyfont=\normalfont,headpunct=:,postheadspace=1em,spacebelow=12pt,spaceabove=2pt]{beweisskizze}
\declaretheoremstyle[headfont=\sffamily\bfseries,bodyfont=\normalfont,headpunct=:,postheadspace=1em,spacebelow=10pt,spaceabove=10pt]{bemerkungen}
\declaretheorem[name=Beweis,numbered=no,style=beweise]{beweis}
\let\proof\relax
\declaretheorem[name=Proof,numbered=no,style=beweise]{proof}
\declaretheorem[name=Sketch of Proof,numbered=no,style=beweisskizze]{sketch}

\declaretheorem[name=Übung,numbered=no,style=bemerkungen]{uebung}
\declaretheorem[name=Erinnerung,numbered=no,style=bemerkungen]{erinnerung}


% -- marginnotes package
\usepackage[fulladjust]{marginnote}
\renewcommand*{\marginfont}{\itshape \footnotesize}
\usepackage{ragged2e}
\renewcommand*{\raggedleftmarginnote}{\RaggedLeft}
\renewcommand*{\raggedrightmarginnote}{\RaggedRight}

% -- todonotes package
\usepackage[textsize=small,obeyDraft]{todonotes}
\LetLtxMacro{\oldtodo}{\todo}
\renewcommand{\todo}[2][]{\tikzexternaldisable\oldtodo[#1]{#2}\tikzexternalenable}
\LetLtxMacro{\oldmissingfigure}{\missingfigure}
\renewcommand{\missingfigure}[2][]{\tikzexternaldisable\oldmissingfigure[{#1}]{#2}\tikzexternalenable}

% -- Fußnoten anpassen
\deffootnote[1.5em]{1.5em}{1.5em}{\textsuperscript{\thefootnotemark}\ }

% -- Figures und Captions
\usepackage{wrapfig}

%--Römische Zahlen
\newcommand{\RM}[1]{\MakeUppercase{\expandafter{\romannumeral #1.\relax}}}
\robustify\RM


% -- Metadaten für die Titelei
\author{Jannes Bantje}


% -- Kopf- und Fußzeilen
\setheadsepline{1pt}[\color{light_gray}]
\pagestyle{scrheadings}
\clearscrheadfoot
\automark{section}
\lehead{\includegraphics[height=0.6 cm,keepaspectratio]{Bilder/Logo_WWU_Muenster_light_gray.pdf}}
\lohead{\color{dark_gray}\normalfont\footnotesize\sffamily\itshape Gruppierung von Trajektorien -- Seminarvortrag von Jannes Bantje}
\rohead{\includegraphics[height=0.6 cm,keepaspectratio]{Bilder/fb10logo_gray.pdf}}

\ofoot[{\color{dark_gray}\sffamily\LARGE \thepage}]{{\color{dark_gray}\sffamily\LARGE \thepage}} %hier wir auch der plain Stil bearbeitet!
\ifoot{ \color{dark_gray} \sffamily\small \leftmark}

% -- Inhaltsverzeichnis und weitere Pakte, die zuletzt geladen werden sollten
\usepackage[tocindentauto]{tocstyle}
\usetocstyle{KOMAlike}
\usepackage{ellipsis}

